\documentclass[11pt]{article}
\usepackage{amsmath,amssymb,amsthm}
\usepackage[margin=1.1in]{geometry}
\usepackage{hyperref}

\newtheorem{proposition}{Proposition}
\newtheorem{remark}{Remark}
\newcommand{\sig}{\sigma}
\newcommand{\D}{\mathcal{D}}

\title{Logistic Regression:\\
A Bernoulli Whose Parameter Is a Function of $X$}
\author{Yuval Lavie}
\date{\today}

\begin{document}
\maketitle

\begin{abstract}
Binary classification, seen from statistics: the label given the features
is \emph{always} a Bernoulli variable --- the only modeling decision is
the form of its parameter function $p(x) = \Pr(Y{=}1 \mid X{=}x)$. The
linear classifier bets on linear log-odds, $p(x) = \sig(w^\top x)$, and
conditional maximum likelihood then reproduces, composed, the two
gradients of the previous notes. New phenomena: the stationarity
equations are transcendental --- the first model in this collection with
\emph{no closed form} --- yet the loss is convex, Newton's method (IRLS)
solves it in a handful of weighted least-squares steps, the fitted
classifier approaches the Bayes floor, and on separable data the MLE
fails to exist. All demonstrated in \texttt{logistic.py}.
\end{abstract}

\section{The conditional-Bernoulli view}

Let $(X, Y)$ be jointly distributed with $Y \in \{0,1\}$. Whatever the
joint $\D$ is, it factors as a marginal over $x$ times a conditional over
$y$ --- and a distribution on $\{0,1\}$ \emph{is} a Bernoulli, so with no
assumption at all:
\begin{equation}
  Y \mid X = x \;\sim\; \mathrm{Bernoulli}\bigl(p(x)\bigr),
  \qquad
  p(x) \;=\; \Pr(Y=1 \mid X=x) \;=\; \mathbb{E}[Y \mid X=x].
  \label{eq:condbern}
\end{equation}
The object of interest is not one random variable but the \emph{family}
of Bernoullis indexed by $x$ --- equivalently the single function
$p:\mathcal{X}\to[0,1]$ (the \emph{regression function}). ``$Z = Y|X$ is
Bernoulli with parameter depending on $X$'' is exactly this, stated for
each slice $X = x$.

\paragraph{Where the assumption actually lives.}
Bernoulli-ness in \eqref{eq:condbern} is forced by the binary support and
costs nothing. The entire model is the choice of a \emph{form} for
$p(x)$. The constant choice $p(x)\equiv p$ is the \texttt{bernoulli/}
experiment; the linear classifier's choice cannot be
$p(x) = w^\top x$ (a line escapes $[0,1]$), so linearity is placed in the
\emph{log-odds}:
\begin{equation}
  \log\frac{p(x)}{1-p(x)} \;=\; w^\top x
  \qquad\Longleftrightarrow\qquad
  p(x) \;=\; \sig(w^\top x) \;=\; \frac{1}{1+e^{-w^\top x}},
  \label{eq:logit}
\end{equation}
with $x$ carrying a fixed leading $1$ so $w$ includes the intercept. In
GLM language: Bernoulli response, logit link. The decision boundary
$\{x : p(x) = \tfrac12\}$ is the hyperplane $w^\top x = 0$ --- the
``linear'' of the linear classifier.

\section{Conditional maximum likelihood}

Given samples $(x_i, y_i)_{i=1}^n$, condition on the observed $x_i$ and
apply \eqref{eq:condbern}: the conditional likelihood and (after the
usual $\log$, $\div n$, negate --- each optimizer-preserving) the average
NLL are
\begin{equation}
  L(w) = \prod_{i=1}^n p(x_i)^{y_i}\bigl(1-p(x_i)\bigr)^{1-y_i},
  \qquad
  \mathcal{L}(w)
  = -\frac1n\sum_{i=1}^n
  \Bigl[y_i\log p(x_i) + (1-y_i)\log\bigl(1-p(x_i)\bigr)\Bigr],
  \label{eq:nll}
\end{equation}
with $p(x_i) = \sig(w^\top x_i)$ --- the \emph{cross-entropy loss},
which by the information-theory note is
$H(\text{data}) + \mathrm{KL}(\text{data}\,\|\,\text{model})$ in
conditional form.

\paragraph{The gradient composes the previous two notes.}
Per sample, write $z_i = w^\top x_i$. The Bernoulli note gives
$\partial \mathcal{L}_i/\partial z_i = \sig(z_i) - y_i$
(the chain rule through $\sig' = \sig(1-\sig)$, everything cancelling);
the linear-regression note gives $\partial z_i/\partial w = x_i$. Chained:
\begin{equation}
  \nabla_w \mathcal{L}_i
  \;=\; \bigl(\sig(w^\top x_i) - y_i\bigr)\, x_i ,
  \qquad
  \nabla_w \mathcal{L}
  \;=\; \frac1n X^\top\bigl(\sig(Xw) - y\bigr)
  \label{eq:grad}
\end{equation}
--- \emph{Bernoulli residual times feature vector}: the same
``prediction minus target, pushed along the features'' shape as least
squares, with the sigmoid interposed.

\section{The closed form dies --- and what survives}

Setting \eqref{eq:grad} to zero:
\begin{equation}
  \sum_{i=1}^n \sig(w^\top x_i)\, x_i \;=\; \sum_{i=1}^n y_i\, x_i .
  \label{eq:stationarity}
\end{equation}
Sigmoids of the unknowns, summed: \emph{transcendental} in $w$. Bernoulli
gave $\hat p = m/n$; least squares gave the normal equations; here
algebra produces nothing --- the first model in this collection where the
theoretical route genuinely does not exist. Not for lack of good
structure:

\begin{proposition}
$\mathcal{L}$ is convex: its Hessian is
\begin{equation}
  \nabla^2 \mathcal{L}(w)
  \;=\; \frac1n X^\top S X \;\succeq\; 0,
  \qquad
  S = \operatorname{diag}\bigl(p_i(1-p_i)\bigr),
  \label{eq:hessian}
\end{equation}
the weights being exactly the conditional Bernoulli variances. Hence
every stationary point is a global minimum (optimization note), and
gradient methods inherit one basin.
\end{proposition}

\paragraph{Newton / IRLS: the statistician's solver.}
Convex + smooth + cheap Hessian invites Newton's method
(optimizers note): $w \leftarrow w - (X^\top S X)^{-1} X^\top(p - y)$
(the $1/n$ cancels). Rearranged, each iteration solves the
\emph{weighted least-squares} system $X^\top S X\, d = X^\top (y - p)$
--- \emph{iteratively reweighted least squares}: logistic regression as a
sequence of linear regressions, reweighted by the current Bernoulli
variances, with quadratic convergence (the script needs 8 steps to
exhaust float64 while GD takes thousands). This is the classical fit
route in statistics packages.

\begin{remark}[When the MLE does not exist: separable data]
If some $w$ classifies the sample perfectly ($w^\top x_i > 0$ exactly
when $y_i = 1$), then scaling $w \to c\,w$ with $c\to\infty$ drives every
$p(x_i)$ to its label and $\mathcal{L} \to 0$ without attaining it: the
infimum is not achieved, $\lVert w\rVert$ diverges, and probabilities
saturate. The script watches $\lVert w\rVert$ grow without bound under GD
on separable data. Cures: $\ell_2$ regularization (adds
$\tfrac{\lambda}{2}\lVert w\rVert^2$, restoring strong convexity and a
finite optimum) or early stopping. Ironically, ``too easy'' data breaks
the estimator.
\end{remark}

\section{The Bayes floor}

Whatever classifier $h$ anyone builds from $x$, its error obeys
\begin{equation}
  \Pr\bigl(h(X) \ne Y\bigr)
  \;\ge\;
  \mathbb{E}\bigl[\min\{p(X),\, 1 - p(X)\}\bigr],
  \label{eq:bayes}
\end{equation}
with equality for the \emph{Bayes classifier}
$h^\ast(x) = \mathbf{1}[p(x) > \tfrac12]$: at each $x$ you cannot beat
predicting the likelier label, and its mistake probability is the smaller
of the two. The floor is a property of the \emph{distribution} (how often
$p(x)$ is near $\tfrac12$), not of any algorithm --- it is the label
noise of the learnability note, seen from the statistics side. Because
our experiment is well-specified (data genuinely logistic), conditional
MLE is consistent for $w^\ast$, and the fitted classifier's error
approaches \eqref{eq:bayes}.

\begin{remark}[Misspecification]
If the true $p(x)$ is \emph{not} of the form \eqref{eq:logit}, the MLE
does not break --- by the information-theory note it converges to the KL
projection of the truth onto the logistic family: the best logistic
approximation, in exactly the divergence that cross-entropy measures.
Well-specified consistency and misspecified KL projection are the same
theorem seen twice.
\end{remark}

\section{Experiment}

$n = 5000$, $x \sim \mathcal{N}(0, I_2)$, $w^\ast = (0.5,\, 2,\, -1.5)$,
labels drawn from \eqref{eq:condbern} (seed 7). Newton/IRLS and GD agree
on the unique optimum ($\hat w \approx w^\ast$ to sampling accuracy);
hand and autograd SGD trajectories coincide to $10^{-10}$; the fitted
classifier's test error matches the Bayes floor
$\mathbb{E}[\min(p, 1-p)]$ to three decimals; and the third panel plots
fitted $\sig(\hat w^\top x)$ against true $p(x)$ on fresh points ---
a diagonal: what was estimated is not a label rule but \emph{the Bernoulli
parameter as a function of $x$}, which is where this note began.

\end{document}
